Los resultados obtenidos muestran que la corriente del diodo en polarización directa sigue el comportamiento exponencial predicho por la ecuación de Shockley, como lo confirma el elevado coeficiente de determinación $R^2=0.998$ del ajuste lineal. Esto indica que la linealización de los datos mediante la representación de $\ln I_d$ en función de $V_d$ fue adecuada dentro del intervalo seleccionado.\\

El valor experimental de la relación ($e/k_B=(1.253\pm0.006)\times10^4\ \mathrm{K/V}$) presentó un error relativo del $7.96 \%$ respecto al valor aceptado. Esta diferencia puede atribuirse a diversas fuentes de incertidumbre, entre ellas la resolución de los multímetros, la variación de la temperatura del diodo durante la medición, la aproximación de considerar constante el factor de idealidad ($n\approx2$) y las simplificaciones realizadas al despreciar la corriente de saturación inversa en la región exponencial. Asimismo, pequeñas variaciones en la selección del intervalo utilizado para el ajuste lineal pueden modificar ligeramente la pendiente obtenida y, en consecuencia, el valor calculado de $e/k_B$. El error se encuentra dentro de un margen razonable para este tipo de prácticas de laboratorio y puede explicarse por las incertidumbres instrumentales, las aproximaciones del modelo y las condiciones experimentales.\\

En polarización inversa se observó que la corriente permaneció muy pequeña dentro del rango de voltajes aplicado, comportamiento consistente con la teoría para un diodo operando por debajo de su voltaje de ruptura. Esto confirma el correcto funcionamiento del dispositivo y la adecuada configuración del montaje experimental.\\
